\section{Anticipated Legal Challenges}
\label{sec:challenges}

Algorithmic redistricting will face legal challenges from political actors who benefit from the current system. This section addresses the most plausible challenges and explains why each fails.

\subsection{Nondelegation}

\textit{Challenge:} Congress cannot delegate the power to draw congressional districts to a computer algorithm. The nondelegation doctrine requires Congress to provide an ``intelligible principle'' to guide any delegation of legislative power.\footnote{See \textit{J.W. Hampton, Jr. \& Co. v.\ United States}, 276 U.S. 394, 409 (1928).} An algorithm is too complex and insufficiently supervised to satisfy this requirement.

\textit{Why this fails:} The nondelegation argument mischaracterizes what algorithmic redistricting does. Congress does not delegate legislative power to the algorithm; it specifies, by statute, the criteria (population equality, compactness, contiguity, no partisan data) and directs executive agencies to implement those criteria using the certified algorithmic method. This is delegation to the executive branch---the Census Bureau---with the algorithm as the Bureau's implementation tool. Congress delegating to an agency with an intelligible principle is constitutionally uncontroversial; the agency using software to implement the delegation raises no additional constitutional question.\footnote{See \textit{Whitman v.\ American Trucking Ass'ns}, 531 U.S. 457 (2001) (upholding delegation with intelligible principle).}

The intelligible principle for algorithmic redistricting is, in any event, far more specific than most delegations the Court has upheld. ``Produce congressional districts that equalize population within $\pm 0.5\%$, minimize boundary perimeter relative to district area, maintain contiguity, and use no data on partisan affiliation or voting history'' is more precise than ``protect public health'' (upholding EPA air quality standards), ``prevent unfair practices'' (upholding FTC authority), or ``prevent excessive rates'' (upholding rate regulation authority).\footnote{See \textit{Yakus v.\ United States}, 321 U.S. 414 (1944); \textit{American Power \& Light Co. v.\ SEC}, 329 U.S. 90 (1946).}

The correct analogy is the Huntington-Hill apportionment statute. Congress mandated that the Census Bureau compute apportionment using a specific mathematical formula. No court has ever suggested this is a nondelegation problem. The formula is the congressional specification; the Bureau's computation is executive implementation. The algorithmic redistricting statute follows the same pattern.

\textit{Nondelegation under} Gundy v.\ United States. Four current justices in \textit{Gundy v.\ United States}, 588 U.S. 128 (2019), indicated readiness to revisit the intelligible-principle doctrine and impose a more demanding standard.\footnote{See \textit{Gundy}, 588 U.S. at 161--62 (Gorsuch, J., dissenting, joined by Roberts, C.J., and Thomas, J.); \textit{Paul v.\ United States}, 140 S. Ct. 342 (2019) (Kavanaugh, J., respecting denial of certiorari) (expressing interest in revisiting the doctrine).} For DIA purposes, the intelligible principle is: ``minimise edge-cut in the census-tract adjacency graph subject to population balance, contiguity, and VRA compliance.'' This is more specific than any contested delegation that has reached the Court---the algorithm itself \textit{is} the regulation, leaving no discretionary gap for an agency to fill through policymaking judgment. The concern that animated the \textit{Gundy} dissenters---that open-ended delegations transfer legislative power from Congress to unaccountable agencies---does not apply where Congress has specified the precise mathematical procedure the agency must execute. A statute that tells the Bureau ``compute districts using these mathematical criteria'' differs categorically from delegations like ``regulate in the public interest'' because the Bureau exercises no policy discretion: it runs a deterministic computation.

\subsection{Elections Clause Limits}

\textit{Challenge:} The Elections Clause grants Congress authority to regulate the ``manner'' of elections, but drawing district boundaries is not the ``manner'' of an election---it is a predicate to elections that establishes the electoral constituencies. Therefore, Congress lacks authority to mandate how states draw districts.

\textit{Why this fails:} This argument was resolved in \textit{Smiley v.\ Holm} (1932). The Court held that the Elections Clause's reference to ``manner'' of elections encompasses the drawing of congressional district boundaries: ``The subject of the time, place and manner of elections ... necessarily includes the authority to require that congressional elections shall be held in particular districts.''\footnote{\textit{Smiley v.\ Holm}, 285 U.S. at 366.} The Court sustained a federal statutory requirement that congressional elections be held in single-member districts as a regulation of the ``manner'' of elections within the scope of Congress's Clause authority.

\textit{Cook v.\ Gralike} (2001), while overturning Missouri's congressional ballot rotation statute as inconsistent with the Elections Clause, confirmed that the Clause grants Congress broad authority to regulate ``how'' elections are held.\footnote{\textit{Cook v.\ Gralike}, 531 U.S. 510 (2001).} A mandate specifying the method by which congressional districts are produced is squarely within this ``how'' authority.

\subsection{VRA Conflict: Algorithm Produces Too Few Majority-Minority Districts}

\textit{Challenge:} The algorithm's baseline mode does not intentionally create majority-minority districts and may produce fewer such districts than enacted plans that have been designed to satisfy VRA requirements. Adopting algorithmic redistricting would violate Section 2.

\textit{Why this fails:} Empirical evidence refutes the factual premise. The algorithm produces more majority-minority districts than enacted plans on average---69 more across the 50-state analysis in \citet{dellaLibera2026vra}---precisely because it does not artificially fragment geographically concentrated minority populations for partisan purposes. Enacted plans in states with Republican-controlled legislatures have frequently drawn minority voters into fewer majority-minority districts than compactness optimization would produce. This is the violation that \textit{Allen v.\ Milligan} addressed; algorithmic redistricting runs in the opposite direction.

Moreover, the VRA does not require intentional creation of majority-minority districts---it prohibits dilution. An algorithm that produces majority-minority districts as a natural consequence of compactness optimization satisfies Section 2 even if it does not target any specific racial percentage. The VRA mode further ensures that wherever a geographically compact minority community exists and satisfies the \textit{Gingles} preconditions, the algorithm will not fragment it.

\subsection{VRA Conflict: Algorithm Produces Racially Gerrymandered Districts}

\textit{Challenge:} The algorithm's VRA mode, which increases edge weights on census tracts with high minority populations, uses race as a factor in drawing districts. Under \textit{Shaw v.\ Reno} and its progeny, this is presumptively unconstitutional racial gerrymandering.

\textit{Why this fails:} \textit{Shaw} targets districts whose shapes are so ``bizarre'' that they are ``unexplainable on grounds other than race.''\footnote{\textit{Shaw v.\ Reno}, 509 U.S. at 644.} The VRA-mode algorithm has a complete non-racial explanation for every boundary it draws: it is minimizing the edge cut of a graph in which certain edges have been assigned higher weights based on the demographic concentration of the census tracts they connect. The result of this optimization is that geographically compact minority communities are preserved---the same outcome that would result from a human mapmaker who chose not to fragment natural communities. The ``bizarreness'' concern does not apply to mathematically compact districts that can be explained by geography.

Furthermore, using race as a factor to satisfy VRA Section 2 is constitutionally permissible provided it is ``narrowly tailored'' to achieving compliance and does not go beyond what compliance requires.\footnote{See \textit{Bush v.\ Vera}, 517 U.S. 952 (1996).} The VRA mode applies edge-weighting to geographic adjacency---it does not target any population percentage threshold but instead prevents fragmentation of geographic minority concentrations. This is a narrow use of demographic information: it preserves what geography created rather than engineering what race requires.

\textit{Race-predominance under} Miller v.\ Johnson. \textit{Shaw} doctrine has evolved through \textit{Miller v.\ Johnson}, 515 U.S. 900 (1995), and \textit{Bethune-Hill v.\ Virginia State Board of Elections}, 580 U.S. 178 (2017), to focus on whether racial considerations \textit{predominated} over traditional redistricting principles, not merely whether district shapes are bizarre. Under the predominant-factor test, a court would ask whether race overrode the algorithm's compactness objective. Under the DIA, no racial data enters the partitioning algorithm. Any VRA adjustment is separately documented under \S\,2(e) and is analytically separable from the base partition: the base partition is produced by the certified algorithm using no demographic data, and the VRA adjustment modifies only specific districts for which a written Gingles justification has been published. Because race modifies the compactness optimization at the margin rather than replacing it, race does not predominate over traditional redistricting criteria---compactness remains the governing principle, with VRA compliance as a documented statutory constraint, in the same way that population equality is a constraint rather than the ``predominant factor'' in every redistricting plan.

\subsection{Communities of Interest}

\textit{Challenge:} Many state constitutions require redistricting bodies to consider communities of interest, including economic, social, and political communities that may not align with geographic compactness. An algorithm that optimizes solely for compactness cannot satisfy this requirement.

\textit{Why this fails:} This challenge conflates what the algorithm produces by default with what the full statutory scheme permits. The model statute provides an adjustment process in which communities of interest may be identified by local advocates, municipal governments, and regional planning bodies. The algorithm's baseline is the starting point, not the final answer. Proposed adjustments that preserve a documented community of interest---a Native American tribal land, a port city's waterfront community, a historically contiguous immigrant neighborhood---are permissible under the statute provided they do not reduce compactness below the statutory threshold and maintain population balance.

More fundamentally, compact districts tend to preserve communities of interest better than non-compact districts. Research by \citet{dellaLibera2026stability} shows that 80\% of algorithmic district boundaries are geographically stable across three census cycles---indicating that the districts reflect underlying community geography rather than arbitrary line-drawing. Districts that track the underlying geography of a state tend to align with communities that have organized around that geography.
